\section{Limitations and conclusion}
\label{sec:limitations}

\subsection{Same-epoch co-movement leaks into finest-scale lagged edges}
\label{sec:lim_leak}

The clearest failure of the model is that same-epoch structure is partly reported
as lagged structure at the finest scale, and the reported edges are lagged self
edges. The benchmark grid places no true edge at the finest scale, so every edge
returned there is a false positive and the rate is directly measurable. Over the
66 fits of the ideal-recovery grid, the local score returns 2.212 spurious
finest-scale edges per fit, 3.244 across the fits that return any at all, and
leaves the finest scale entirely clean in 21 of them. Over the 54 fits of the
algorithm comparison the same leak is a saturation of 0.293 of the one-parent
budget per label. That is far better than every classical entry, five of which
use the entire budget for every label in every fit and three of which reach
0.770 to 0.967, and it is not zero.

Two qualifications belong with those rates. The first is that the grid was
reshuffled between this run and the one reported before it, which moved true
edges onto the coarsest scale, so the earlier finest-scale rate of 1.106 edges
per fit and the present 2.212 are not a like-for-like pair. What can be said is
that the direction is against the model. The second is that the finest scale is
now known to be unable to carry a findable edge at all. Run at a nonzero
in-degree on every scale, the screen of \cref{sec:benchmark_identifiability}
puts the median gain of a true finest-scale edge below the floor it has to clear
in all three families, by 3.40, 9.46 and 6.27 nats, with three of twenty-one
scenarios clearing, because the finest scale is a binary event-start process and
a lagged edge there can only move which minute of a block begins an episode. A
method that returns an edge there is reporting structure the data does not
contain, and no measurement on this grid separates that from correctly finding a
finest-scale edge, since none can be placed.

Four candidate causes were tested. Three were excluded by measurement and the
fourth is what survives. The first candidate was the latent factor, on the
reading that the leak is co-movement a factor would otherwise absorb. It is
excluded, because the scenario families that carry a factor and force a single
regime leak at 0.289 in saturation, essentially the 0.283 of the family that
carries no latent state at all. The second candidate was the regime, whose
persistence autocorrelates every label it drives. The same measurement excludes
it, since the families that carry a regime and no factor leak at 0.296, the same
rate again, and the stationary family, which has neither a regime nor a factor,
still leaks at 0.283 over its 12 fits, 6 of which return a finest-scale edge. On
the grid used before this one the regime carried a visible share of the
difference between the families and was excluded only as the sole cause; on this
one the families do not differ, so there is no difference left for it to carry.
The third candidate was the finest-scale cap that constrains counts to be binary,
on the reading that the allocation likelihood does not model the redistribution
the cap forces. It is excluded, because a probe with the cap removed leaks more
rather than less.

What survives all three exclusions is within-block competition for a fixed parent
total. The children of a parent epoch share a total that does not change, so a
child whose sibling has already taken counts has fewer left to take, and a lagged
feature of a sibling in the same block is a proxy for how much of the block total
is already used up. Three properties of the measurement identify this mechanism
rather than any of the excluded ones. The effect is present within a block and
absent across block boundaries; the coefficient the search fits for the spurious
edge is negative, which is the sign competition predicts and the sign a genuine
excitatory dependence does not; and the effect survives in the family with no
latent state at all, because the block total is fixed there too. Adding an
explicit depletion covariate to the allocation score, an exclusive cumulative sum
of the counts already taken by earlier children of the same block, captures a
large amount of real structure and improves held-out fit, and it does not remove
the finest-scale false positives of the benchmark. The gains those surviving
edges carry sit an order of magnitude above anything a block-preserving surrogate
of the same feature earns, which places their source in the part of the
finest-scale regime path the posterior cannot identify rather than in
competition.

The refactorization that looks like the obvious closure is not one. The
Dirichlet-multinomial log-likelihood of one label in one block, with parent total
\(N\), child counts \(x_t\) summing to \(N\), allocation probabilities
\(\omega_t\) summing to one over \(t \in \child_k(c)\), and concentration
\(\kappa\), can be written as a sequential product of beta-binomial kernels, one
per child given its realized predecessors, which is exactly the form an exact
per-child normalizer would need. Ordering the children arbitrarily and setting
\(a_t = \kappa \omega_t\),
\(b_t = \kappa \sum_{s>t} \omega_s\)
and \(N_t = N - \sum_{s<t} x_s\), the identity is
\begin{align}
&\sum_{t \in \child_k(c)}
\Bigl[
\log \Gamma\!\left(x_t + \kappa\omega_t\right)
- \log \Gamma\!\left(\kappa\omega_t\right)
\Bigr]
+ \log \Gamma(\kappa)
- \log \Gamma(N + \kappa)
\nonumber\\
&\qquad =
\sum_{t \in \child_k(c)}
\Bigl[
\log B\!\left(x_t + a_t,\, N_t - x_t + b_t\right)
- \log B\!\left(a_t, b_t\right)
\Bigr],
\label{eq:lim_betabinomial}
\end{align}
where the last child in the chosen order has \(b_t = 0\) and \(x_t = N_t\), so
the two divergent beta terms it contributes cancel and its summand is zero, which
is the statement that the last child takes whatever the block has left. The
identity holds for every \(\omega\), including an \(\omega\) that depends on the
data, and the combinatorial factors the two sides omit are equal to each other,
since \(\prod_t \binom{N_t}{x_t} = N! / \prod_t x_t!\). Both sides are therefore
the same function of the parameters, with the same gradient, the same maximizer
and the same edge gain. A direct evaluation of the two sides agrees to
floating-point rounding. Refactorizing the block changes no
score, so it is not a route to an exact normalizer, and an earlier claim in this
project that it was such a route is withdrawn.

\subsection{The allocation normalizer is evaluated under a mean field}
\label{sec:lim_meanfield}

The allocation probabilities of a block couple all of its children through a
single normalizer, and the implementation evaluates that normalizer by scoring
each child at its own per-regime score while holding its block-mates at the
regime-averaged score. The consequence is precise. The parameter step accepts a
block when the objective computed at the regime posterior it was given does not
fall, and the next latent-state inference step recomputes that objective at the
new posterior. The two are not the same function, so a block that raised the
first can lower the second, and the second is what the outer loop monitors.

The remedy shipped here is an outer guard rather than an exact normalizer. A
sweep that ends below the best sweep so far is rejected, the state reverts to the
best sweep, and the sweep is retried with a halved parameter learning rate up to
a bounded number of times, after which the fit stops at its best state. The guard
works as specified. Over the 1056 fits of the reported runs that execute more
than one sweep, the recorded objective decreases in none, against a decrease in
81.1 percent of the corresponding fits of the run made before the guard existed.
The residual is contained and not removed, and the count of rejected sweeps is
the measure of how much of it there is. The guard rejected at least one sweep in
826 of the 1176 reported fits, at most three in any fit. Removing the residual
needs a normalizer that is exact under the variational family, and the exact
expectation of the block log-normalizer under the joint posterior over the
children's regimes has a number of terms exponential in the block size, so it is
not available at the branching factors used here.

\subsection{Recall on external data}
\label{sec:lim_recall}

On released Bayesian networks tiled into two-slice temporal networks, hill
climbing under the Bayesian information criterion still recovers more edges than
the local score does. Most of the gap the first version of this experiment
showed was the protocol rather than the criterion, and the part that survives
correcting the protocol is smaller and has a structural explanation.

Three protocol choices were wrong and all three were ours. Two regimes were
being fitted to data generated with none, so the regime chain was free to absorb
variation the emission owns. The edge penalty was inherited from the synthetic
grid, where a persistent latent state makes a non-edge expensive to reject,
while a tiled network has no such null level. And the block length of the
imposed hierarchy was set to ten, in the flat part of no curve. Selecting the
regime count, the factor dimension and the penalty on held-out log evidence, and
declaring the block length as the square root of the series length, takes $F_1$
on five tiled networks from $0.532$ to $0.806$, precision from $0.750$ to
$0.869$ and recall from $0.414$ to $0.762$. Two of the five are then recovered
as well as hill climbing recovers them, exactly $F_1 = 0.933$ on asia and
exactly $1.000$ on earthquake.

What remains is concentrated rather than uniform. On child, the largest network
scored, the local score reaches $F_1 = 0.870$ against $0.933$ at the higher
precision of the two, $0.958$ against $0.933$. On sachs it reaches $0.812$
against $0.980$; on cancer it reaches $0.413$ against $0.905$.

Averaged over the five networks and three replicates the deficit is about eight
true edges, spread over three networks: $3.33$ on child, $2.67$ on sachs and
$2.00$ on cancer, with asia and earthquake exactly tied. On the single
replicate we examined in detail the deficit is three edges, and we report those
three rather than a mechanism we cannot support. On sachs the two are
\texttt{PKA} to \texttt{Mek} and \texttt{Erk} to \texttt{Akt}. The first is a
budget substitution: the in-degree was saturated at three of three and the slot
went to \texttt{Plcg}, whose conditional gain is $-13.88$. The second is not:
the in-degree was two of three with a slot free, and the score declined to add
an edge whose conditional gain is $+17.36$. On cancer the score returns no
parents at all for \texttt{Cancer}, although \texttt{Smoker} to \texttt{Cancer}
has a sole-parent gain of $+20.91$ and an odds ratio of $15$. A positively
scoring edge that is not taken with room to take it is a defect in the local
score or in what the latent posterior does to it, and we state it as an open
question rather than explain it.

We had proposed that the finest-scale term, being a within-block allocation
conditioned on the block total, cannot see a parent effect that moves the total.
That explanation is wrong and we withdraw it. On the three edges at issue the
within-block association is the whole association: the ratio of within-block to
raw association averages $1.007$ for recovered edges against $0.999$ for missed
on sachs, and $1.003$ against $0.989$ on cancer; conditioning exactly on block
totals retains $99.4$ against $98.7$ percent of the pooled association on sachs
and $97.7$ against $93.5$ percent on cancer; and between-block variation is
$1.1$ to $2.0$ percent of each target's fine-scale variance against $1.6$
percent expected from sampling alone. There is no between-block variation for
an effect to hide in.

Two candidate explanations were tested and rejected. The search is not the
cause, since exhaustive per-target enumeration and a genetic search return the
same parent sets as the greedy one on every cell of both benchmarks. And
error-rate-controlling selection is not the remedy, since subsample stability
raises precision to $1.000$ but costs enough recall to lower $F_1$, and held-out
evidence prefers the single fit.

\subsection{Parameter recovery fails the ladder from either start}
\label{sec:lim_params}

Parameter recovery is a consistency claim, so the test that was stated before
the experiment ran is a sample-size ladder rather than any single value, and the
criterion was that the error must fall at every step. The ladder is set from the
coarsest scale rather than the finest, so its three rungs give that scale $168$,
$336$ and $672$ epochs, and the middle rung is the horizon the structural
benchmark itself runs at. Neither fitted start meets the criterion. From a
ground-truth start the mean parameter error runs $0.2899$, $0.3494$ and
$0.3208$; from a data-driven start $0.4261$, $0.4533$ and $0.4678$, rising at
every rung.

The failure is confined to one generating family. Family III falls at every
rung, $0.2956$, $0.2932$ and $0.2588$, and meets the criterion outright; family
I is close to flat; family II rises at every rung, $0.6461$, $0.7933$ and
$0.8445$, and sets the level. An error that rises with the sample size is a fit
converging to a different solution and describing that better, rather than a
slow one, and family II is the regime-carrying family whose finer latent paths
are unidentified even with the generating parameters installed.

One part of the model's behavior at the coarsest scale has a mechanism and a
correction. A level initialized from the observed mean already contains the
contribution of the lagged parents, and the parameter step refits that level
against a graph with no edges in it before the structure search runs, so the
level absorbs exactly what a true lagged edge would explain and the search meets
a model that already accounts for the data. Correcting both, by initializing the
level net of what its own lag-1 values explain and by searching before the level
is refitted, takes coarsest-scale recovery from $0.00$, $0.00$ and $0.67$ lagged
keys per fit to $4.67$, $4.44$ and $4.67$ against a true $4.44$. Measured against
the same fits without the correction, parameter error improves in seven of the
nine family-by-rung cells and worsens in one, by $0.044$. Family II's rise is
present in both, so it is not caused by the correction.

The practical consequence is unchanged. Fitting an observed log supports
prediction, exact reconciliation and stability selection, and does not yet
support a claim about recovered parameters, because the criterion set in advance
is not met from the only start a real log admits.

\subsection{What no external benchmark can test}
\label{sec:lim_scope}

The scope limit follows from \cref{sec:related_benchmarks} and cannot be
engineered away within this paper. No released benchmark combines event counts,
exact cross-scale aggregation and a latent regime, so the property the model is
built around cannot be scored against an external ground truth. The tiled
networks test the structure learner, the local score and the table estimator on
flat discrete data, and the two-scale hierarchy imposed on them by summing blocks
of epochs is an artifact of the experiment rather than a property of the data.
Every multiscale number in this work therefore comes from a generator written for
it, and the defense of that generator is the identifiability screen rather than
an appeal to realism.

What a real log can test is prediction, exact reconciliation, the stability of a
selected edge across subsamples, and the shift confirmation used inside the
structure step, since none of these needs a known graph. The self-compatibility
check of \citet{faller2024selfcompatibility} adds a further falsification test
that requires no ground truth, and it is necessary rather than sufficient.
Recovery of a graph, of a regime path or of a parameter cannot be measured on
real data at all, and no claim of that kind is made here from real data.

\subsection{What the observational arm does not settle}
\label{sec:lim_arm1}

The observational arm of \cref{sec:results:arm1} is the one place a real log is
put to the model, and four of its results are bounded in ways a reader has to
carry.

The grid's reach is a property of the roster. Of $228$ cells, $46$ are put to a
test and $127$ of the rest are descriptive because the declared tolerance is
inside the precision $51$ participants over eight weeks can resolve. Three of
the four declared negative controls fail to break at least one row they are
declared against. A control that cannot break a row is measuring what the grid
can reach on this panel, so $46$ tested cells is the sensitivity of the
comparison and not a property of the fitted model. A larger roster would move
that number and might move the verdicts with it.

Five of the $11$ tested rows fail their own replicate noise, which is why the
row-level reading of \cref{sec:results:arm1} is reported net of the positive
control and the raw count of $12$ failing cells of $46$ is not the figure to
read. A tolerance narrower than the spread between two draws of the same fitted
model grades the draw's variance and says nothing about the fit. The same
reading bounds the comparison against the empty graph. The fitted leg fails two
cells fewer of $46$, and the single row that separates the two is one of the
five that fail their own noise, so on the six rows that carry evidence the
fitted graph and an intercept with a participant level pass and fail exactly the
same ones. The grid is not what separates them on this roster.

Two of the six labels are drawn at rates the panel cannot pin down.
\texttt{mindfulness} occurs once per $40$ observed participant-days and is
drawn at $1.94$ times its rate; \texttt{survey\_response} is a count-only node
whose duration block is uninformative and is drawn at $0.53$. The drawn logs
also leave $5.3$ percent of observed days with no episode against the panel's
$17.0$ percent, a marginal quantity no row of the grid reads. These are
reported as measured and none of them is repaired by a correction inside the
model as it stands.

No edge between two distinct behavior labels is bought at any scale, so the
graph this arm reports is a set of clock and scheduling effects and a set of
per-label persistences. That is a negative result about this panel and it is
reported as one. Nothing in the declared anchor set adjudicates it, since every
declared positive is the scheduled notification into the survey node, and the
arm supplies no instrument that separates a between-label relation this cohort
does not carry from one eight weeks of $51$ participants cannot price above the
null.

The graph the arm reports carries two edges whose coefficients cannot be
interpreted. Two lag-two self edges into \texttt{nap} at the hour scale reach
standardized effects of $19.2$ and $-145.0$ on columns that clear the support
floor and carry far too little information to determine what they are given.
They are reselected on one of $16$ half-panels each, so the stability rung
removes them, and the support floor is the wrong instrument for this case. The
invariance rung rejects the model for \texttt{gym}, on the one accepted parent
that target carries at the hour scale, so it is a weak test and it is reported
with its width printed. Neither result is repaired here, and both are reported
beside the graph they qualify.

\subsection{Conclusion}
\label{sec:conclusion}

This work establishes five things.

A single model can meet the four requirements of
\cref{sec:intro_requirements} at once. Exact aggregation between a fine count and
its coarse summary holds by construction rather than by repair, and the measured
reconciliation error is zero for counts and durations in every reported fit,
including a fit on an observed log. Lagged dependence is learned per scale under
a penalized local score rather than assumed. Persistent drift and same-epoch
co-movement are carried by a regime and a factor respectively, so that
co-movement is not reported as a directed dependence. And the finest scale maps
back deterministically to a dated event trace.

Scoring a candidate parent set against a model that already contains the latent
variables changes what the search finds. On a benchmark screened so that its
scenarios contain what the metrics claim to measure, over 540 fits spanning 18
scenarios at three replicates and ten methods, the latent-conditional score
reaches \(F_1 = 0.690\) at the ten-minute scale and \(0.892\) at the hourly
scale, the two that carry true edges, against 0.109 to 0.251 and 0.021 to 0.346
for eight published learners, and it is the only entry whose total structural
Hamming distance, 5.963, beats the 10.981 of reporting nothing. Every classical
entry lands between 1.98 and 2.92 times that reference. The mechanism behind the
gap is a confound any latent-free criterion shares, and recall split by edge kind
states it directly rather than leaving it to be inferred. At the ten-minute scale
seven of the eight latent-free entries recover between 0.679 and 1.000 of the
true lagged self edges, and none of the eight recovers more than 0.245 of the
cross-label ones, while the latent-conditional score recovers 0.919 and 0.523 and
is the only entry above a quarter on cross edges.

On external ground truth the ordering depends on how the two model classes are
configured, which is itself a finding about the comparison rather than about
either criterion. A latent-variable model carries free quantities that a
latent-free learner does not, and holding them at declared values while the
classical learners choose their own penalty through their own score measures
the declarations. Selecting the regime count, the factor dimension and the edge
penalty on held-out log evidence from the training split raises $F_1$ on the
tiled networks from $0.532$ to $0.806$ when each dataset selects for itself,
and to $0.869$ at a precision of $0.983$ when one penalty is chosen for the
benchmark. On the synthetic grid the same rule does not select the maximizing
penalty. It picks $0.5$, where $F_1$ at the ten-minute scale is $0.610$, while
that $F_1$ peaks at the strictest setting of $2.0$ with $0.691$, which is also
the best setting on precision, $0.900$, and on total distance, $6.82$. Held-out
likelihood does not order structure where a persistent latent state
makes a spurious edge a good predictor, which is a limitation of the selector
rather than of the score it selects for. The frontier was once scored on the
ten-minute scale alone, which left open whether the disagreement was an artifact
of reading a two-scale truth through a one-scale metric. Rescoring at the
coarsest scale and on the mean closes that: both disagree with the selector as
well, at $1.0$ and $2.0$ against the selected $0.5$. What the wider reading adds
is that the two truth-bearing scales do not agree with each other, so no scalar
penalty is best at both. Hill climbing under BIC still leads on the
tiled networks, by a margin in recall, so the statement is that the gap is
smaller and one-sided rather than that the two benchmarks contradict each other.

A model fitted to a real log draws a log a panel of non-stationarity tests
cannot easily tell from it. On $51$ LifeSnaps participants over eight complete
calendar weeks each, the fitted graph draws four of six labels within $6$
percent of their real episode volume at exactly zero reconciliation error and
four times closer than an intercept with a participant level draws them; net of
a positive control that names the rows failing their own replicate noise, it
passes four of the six grid rows that carry evidence. On those six rows the
empty graph passes and fails exactly the same ones, so the grid at this roster
size bounds what it can resolve and the volume comparison is what separates the
two. Changepoint agreement between the two
logs is monotone in the calendar scale for all six labels, from $0.95$ to
$0.96$ at the week scale for three of them down to $0.05$ at the hour scale for
the one count-only label, which places the departure at the within-day shape
and nowhere coarser. The same panel confirms
\cref{cor:model_calendar_home} in both of its directions: offering the calendar
nodes at their home scale is worth $2295.64$ nats over withholding them and
offering them at every scale costs $33.90$ nats against offering them at their
home scale. On that log every cross-label edge the search buys sits at the hour
scale and every edge it buys at the day and week scales reads a label's own
past, which is the block decomposition of
\cref{prop:model_block_invariance} appearing as a property of a fit.

Finally, the failures are located rather than described. The finest-scale leak
has three candidate causes excluded by measurement and one algebraic dead end
disproven; the mean-field normalizer is contained by a guard whose residual is
counted rather than assumed away; the residual external recall gap is
quantified after selection rather than argued away; and parameter recovery is
reported as failing a criterion stated in advance from either fitted start, with
its candidate cause named as a hypothesis to test rather than as a conclusion.
On the observed log the search buys no edge between two distinct behavior
labels at any scale, and the declared anchor set contains no such relation to
adjudicate it, so that count is reported with the two readings it admits and
with the statement that this arm carries no instrument to choose between them.
